\documentclass[11pt]{article}
\usepackage[a4paper,margin=2.3cm]{geometry}
\usepackage[T1]{fontenc}
\usepackage{amsmath,amssymb}
\usepackage[hidelinks]{hyperref}

\title{Simplex integration by recursive convolution}
\author{}
\date{}

\begin{document}
\maketitle
\vspace{-1.5em}

We start from a hypercube integral, restrict it to an ordered simplex
using permutation symmetry, and reduce the simplex integral to successive
one-dimensional convolutions. Physical prefactors, such as powers of the
interaction or factors of $i$, are omitted; the combinatorial factor
$1/{n!}$ is kept explicitly until the change of domain.

\section{Original integration on the hypercube}

Start from the integrand $\widetilde Q_n(u_1,\ldots,u_n)$ and the
normalized integral
\begin{equation}
    \boxed{Q_n(t)=\frac{1}{n!}
    \int_{\mathcal H_t}du_1\cdots du_n\,
    \left[\prod_{k=1}^n\lambda(u_k)\right]
    \widetilde Q_n(u_1,\ldots,u_n),
    \qquad \mathcal H_t={[0,t]}^n.}
    \label{eq:hypercube_integral}
\end{equation}
All variables are integrated independently over the same interval.
The function $\lambda$ describes the switching on of the interaction;
$\lambda=1$ on the interval gives the unweighted case.
For identical interaction vertices, relabelling the vertices does not
change the complete time-ordered diagrammatic expression, with its
fermionic signs already included. After summing any associated internal
labels, this gives permutation symmetry in the integration times.
We use this property of the full integrand:
\begin{equation}
    \widetilde Q_n(u_{\pi(1)},\ldots,u_{\pi(n)})
    =\widetilde Q_n(u_1,\ldots,u_n)
    \qquad\text{for every permutation }\pi.
    \label{eq:symmetry}
\end{equation}
The activation product is also permutation symmetric, so the complete
weighted integrand has this property.

\paragraph{Why not factorize directly on the original hypercube?}
Permutation symmetry does not imply smoothness or a low-rank tensor
representation. Time-ordered propagators contain Heaviside factors such
as $\theta(u_i-u_j)$, and the resulting integrand can have cusps when two
interaction times coincide. The hyperplanes $u_i=u_j$ separate the
hypercube into ${n!}$ ordering sectors. A direct tensor approximation in
the $u_k$ coordinates must resolve these internal ordering boundaries,
which can require large ranks even when the function is simpler within
each sector.

The purpose of the following change of variables is therefore not merely
to replace one integration domain by another. We first select one
ordering sector using symmetry, then express the integrand in successive
time gaps. Coincident adjacent times become boundaries at zero gap,
rather than diagonal boundaries inside the domain. This motivates
learning the integrand in gap coordinates instead of the original time
coordinates. The gap-coordinate function may be approximated on a
hypercube, while the physical simplex constraint is imposed during
integration; it is not inserted as a discontinuous mask into the function
to be factorized.

\section{Passage from the hypercube to the simplex}

Define the ordered simplex
\begin{equation}
    \mathcal S_u(t)=\{(u_1,\ldots,u_n):
        0\leq u_n\leq\cdots\leq u_1\leq t\}.
    \label{eq:ordered_simplex}
\end{equation}
The hypercube splits into ${n!}$ ordering sectors, one for each permutation
of the variables. Their interiors are disjoint, and their common
boundaries $u_i=u_j$ have zero measure for an ordinary integrable function.
By symmetry, each sector has the same integral. The hypercube integral
is therefore ${n!}$ times the integral over $\mathcal S_u(t)$, cancelling
the original factor $1/{n!}$:
\begin{equation}
    \boxed{\begin{aligned}
    Q_n(t)&=\int_{\mathcal S_u(t)}du_1\cdots du_n\,
        \left[\prod_{k=1}^n\lambda(u_k)\right]
        \widetilde Q_n(u_1,\ldots,u_n)\\
    &=\int_0^t du_1\int_0^{u_1}du_2\cdots
      \int_0^{u_{n-1}}du_n\,
      \left[\prod_{k=1}^n\lambda(u_k)\right]
      \widetilde Q_n(u_1,\ldots,u_n).
    \end{aligned}}
    \label{eq:ordered_integral}
\end{equation}
No additional division by ${n!}$ is needed after this step.
If the starting hypercube integral has no $1/{n!}$ prefactor, its simplex
expression instead retains an overall factor ${n!}$. Without permutation
symmetry, the different ordering sectors must be summed explicitly.
For a common interval $[a,b]$, the same construction applies after
shifting $u_k=a+s_k$, setting $t=b-a$, and replacing $\lambda(s)$
by $\lambda(a+s)$.

\paragraph{Successive time gaps.}
Introduce the successive time gaps
\begin{equation}
    v_1=t-u_1,\qquad v_k=u_{k-1}-u_k\quad(k\geq2),
    \qquad u_k=t-\sum_{j=1}^k v_j.
    \label{eq:gaps}
\end{equation}
The absolute Jacobian is one. The ordered simplex becomes
\begin{equation}
    \mathcal S_v(t)=\left\{(v_1,\ldots,v_n):
        v_k\geq0,\quad\sum_{k=1}^n v_k\leq t\right\}.
    \label{eq:gap_simplex}
\end{equation}
For the chosen upper limit $t$, define the integrand in gap coordinates by
\begin{equation}
    f_n(v_1,\ldots,v_n)=
    \widetilde Q_n\left(t-v_1,\ t-v_1-v_2,\ \ldots,
                       t-\sum_{j=1}^n v_j\right).
    \label{eq:gap_integrand}
\end{equation}
The same integral can then be written
\begin{equation}
    Q_n(t)=\int_{\mathcal S_v(t)}dv_1\cdots dv_n\,
      \left[\prod_{k=1}^n
      \lambda\left(t-\sum_{j=1}^k v_j\right)\right]
      f_n(v_1,\ldots,v_n).
    \label{eq:gap_integral}
\end{equation}
The two domains are the same simplex in different coordinates.
Adjacent coincident times $u_{k-1}=u_k$ become boundary faces $v_k=0$;
the ordering boundaries no longer cut through the interior of the domain.
In the ascending delay convention $w_k=t-u_k$, one instead has
$0\leq w_1\leq\cdots\leq w_n\leq t$, with
$v_1=w_1$ and $v_k=w_k-w_{k-1}$.
Thus, a function of the individual delays is evaluated in gap coordinates
at $(w_1,\ldots,w_n)=(v_1,v_1+v_2,\ldots,\sum_j v_j)$.

\section{Tensor-train reduction to one-dimensional integrals}

To obtain the full time dependence from one construction, assume that
the gap-coordinate integrand admits a common tensor-train approximation
throughout the time window $0\leq t\leq T_M$, with fixed kernels and
pivot matrices:
\begin{equation}
    f_n(v_1,\ldots,v_n)
    \simeq T_1(v_1)P_1^{-1}T_2(v_2)\cdots
           P_{n-1}^{-1}T_n(v_n).
    \label{eq:tt}
\end{equation}
Here $T_k(v)$ is an $r_{k-1}\times r_k$ matrix, $r_0=r_n=1$,
and $P_k$ is an invertible $r_k\times r_k$ cross-interpolation pivot
matrix. These are ordinary matrix products, without conjugation.
The pivot inverses may be absorbed into the kernels by defining
\begin{equation}
    A_k(v)=T_k(v)P_k^{-1}\quad(k<n),\qquad A_n(v)=T_n(v).
\end{equation}

Integrate the innermost variable $u_n$ first, then $u_{n-1}$, and so on.
For $k<n$, the recursion is
\begin{equation}
    \boxed{\begin{aligned}
        \Psi_k(x)&=\int_0^x dy\,
            \lambda(y)T_k(x-y)P_k^{-1}\Psi_{k+1}(y),
            &&k=n-1,\ldots,1,\\
        \widehat Q_n(t)&=\Psi_1(t)\simeq Q_n(t).
    \end{aligned}}
    \label{eq:recursion}
\end{equation}
The vector $\Psi_k(x)$ has $r_{k-1}$ components; the last result
$\Psi_1(t)$ is a scalar. In particular,
\begin{equation}
    \Psi_n(x)=\int_0^x\!\mathrm{d}y\,\lambda(y)T_n(x-y).
\end{equation}
To derive the recursion, observe that $u_n$ occurs only in
$\lambda(u_n)T_n(u_{n-1}-u_n)$. Integrating this factor over
$0\leq u_n\leq u_{n-1}$ gives $\Psi_n(u_{n-1})$.
Now suppose that eliminating $u_{k+1},\ldots,u_n$ has produced
$\Psi_{k+1}(u_k)$. The remaining factor involving $u_k$ is
\begin{equation*}
    \lambda(u_k)T_k(u_{k-1}-u_k)P_k^{-1}\Psi_{k+1}(u_k).
\end{equation*}
Its integral over $0\leq u_k\leq u_{k-1}$ is precisely
$\Psi_k(u_{k-1})$. Repeating this step reaches $u_0=t$ and leaves
$\Psi_1(t)$, proving the reduction while preserving the matrix order.

With the causal matrix-vector convolution
\begin{equation}
    (A*B)(x)=\int_0^x\!\mathrm{d}y\,A(x-y)B(y),
\end{equation}
each step is simply $\Psi_k=A_k*(\lambda\Psi_{k+1})$,
with $\Psi_{n+1}(x)=1$ for the base case.
The activation multiplies the current vector \emph{pointwise before each
convolution}; it is not applied only once at the end.
For example, at order two,
\begin{equation}
    \widehat Q_2(t)=\int_0^t\!\mathrm{d}u_1\int_0^{u_1}\!\mathrm{d}u_2\,
       \lambda(u_1)\lambda(u_2)
       T_1(t-u_1)P_1^{-1}T_2(u_1-u_2).
\end{equation}

\paragraph{Final integration and the complete time dependence.}
For $n\geq2$, the last remaining integration is
\begin{equation}
    \boxed{\widehat Q_n(t)=
    \int_0^t\!\mathrm{d}u_1\,\lambda(u_1)T_1(t-u_1)P_1^{-1}\Psi_2(u_1)
    =\Psi_1(t).}
    \label{eq:correct_final}
\end{equation}
The important point is that each $\Psi_k$ is constructed as a
\emph{function} on $[0,T_M]$, not as a scalar for one chosen observation
time. Once $\Psi_2$ has been constructed, the last convolution constructs
$\Psi_1$ over the entire time window. Consequently,
\begin{equation}
    \boxed{\widehat Q_n(t)=\Psi_1(t)
        \qquad\text{for every }0\leq t\leq T_M.}
\end{equation}
The full time dependence is thus obtained in the same final
function-construction step: no new integrand fit and no repetition of
the preceding integrations are required for each time point. Changing
$t$ only requires evaluating the resulting function. This is the key
advantage over calculating a separate multidimensional integral at each
observation time. The construction cost still depends on the chosen
resolution and ranks; explicitly exporting many values also has an
output cost.

This reuse requires the common gap-coordinate representation assumed
above. For a generic starting integrand, changing coordinates alone does
not guarantee kernels independent of the chosen upper limit.
For $n=1$, the base convolution already constructs the whole output
function, with no pivot matrix.

\section{How to perform the integration with quantics}

\subsection{Binary grid and the two kinds of tensor indices}

Choose a common grid for the gaps and the intermediate time variable:
\begin{equation}
    N=2^B,\qquad h=\frac{T_M}{N},\qquad x_m=mh,
    \qquad m=\sum_{\ell=0}^{B-1}2^\ell b_\ell,
    \quad b_\ell\in\{0,1\}.
\end{equation}
The stored indices are $m=0,\ldots,N-1$. Quantics replaces each sampled
one-dimensional function by an MPS with $B$ binary physical indices,
ordered here from the least to the most significant bit. In particular,
represent the arrays
\begin{equation}
    \Lambda[m]=\lambda(mh),\qquad
    \mathsf T^{(k)}_{ab}[q]={\bigl(T_k(qh)\bigr)}_{ab},
    \qquad \Psi_{k,a}[m]\simeq{\bigl(\Psi_k(mh)\bigr)}_a
\end{equation}
as scalar MPSs, one for each component. The indices $a,b$ are the
\emph{variable-level} train indices of Section~3; they are not binary
physical indices or quantics bond indices. This componentwise construction
is one direct way to retain the matrix structure of the original train
without constructing dense arrays of length $N$.

\subsection{The recursion becomes MPS products, sums, and convolutions}

Use the raw, nonperiodic discrete convolution
\begin{equation}
    \operatorname{conv}(A,D)[m]=\sum_{q=0}^m A[q]D[m-q],
    \label{eq:raw_quantics_conv}
\end{equation}
which contains no factor $h$. To choose a definite quadrature, take the
left-rectangle rule in the integration variable $y$ of Section~3.
It uses $y=0,h,\ldots,(m-1)h$, or equivalently positive lag indices
$q=m-y/h=1,\ldots,m$. Thus define kernels with their zero-lag entry removed:
\begin{equation}
    \mathsf K^{(k)}_{ab}[q]=(1-\delta_{q,0})\mathsf T^{(k)}_{ab}[q].
    \label{eq:positive_lag_kernel}
\end{equation}
In MPS form this can be done by subtracting
$\mathsf T^{(k)}_{ab}[0]\,\delta_{q,0}$; the delta at the all-zero bit
string is a bond-dimension-one MPS\@. This removes a quadrature endpoint,
not the value of the original continuous kernel.

Starting at $k=n$ and proceeding downwards, the discrete recursion is
\begin{equation}
    \boxed{\begin{aligned}
    \Psi_{n,a}
        &\simeq h\,\operatorname{conv}(\mathsf K^{(n)}_{a1},\Lambda),\\
    D_{k,b}
        &=\Lambda\odot\sum_{c=1}^{r_k}{(P_k^{-1})}_{bc}\Psi_{k+1,c},
        &&k<n,\\
    \Psi_{k,a}
        &\simeq h\sum_{b=1}^{r_k}
            \operatorname{conv}(\mathsf K^{(k)}_{ab},D_{k,b}),
        &&k<n.
    \end{aligned}}
    \label{eq:quantics_recursion}
\end{equation}
Here $\odot$ denotes the pointwise product of two MPSs representing
functions on the same grid. Mixing by $P_k^{-1}$ uses scalar multiples
and sums of component MPSs. Each convolution is a scalar MPS convolution;
the component sums restore the matrix-vector contraction.
Truncation may be applied after these operations to control the quantics
bond dimensions. No complex conjugation is involved.

Each recursion level contributes one factor $h$, so the final result
contains $h^n$ already. Evaluate the scalar MPS $\Psi_{1,1}$ at the bits
of $m$ to obtain the approximation to $Q_n(mh)$; do not multiply by a
second volume factor or divide by ${n!}$ again. The point $T_M$ itself is
not stored, so evaluating there requires an endpoint convention rather
than silently interpreting $m=N-1$ as $m=N$.

\subsection{What this means for a TTN fitted directly in the gaps}

The componentwise procedure above assumes the variable-level train
$T_1P_1^{-1}\cdots T_n$. A general TTN does not directly provide these
matrix kernels. One possibility is to convert it to such an ordered train,
but the variable-level ranks can grow. An alternative is to keep the TTN
and attach quantics arithmetic constraints to its branches.

To see precisely which constraints are required, let $j_k$ index the
$k$th gap and introduce a remaining-time index after each gap:
\begin{equation}
    s_0=m,\qquad s_k=s_{k-1}-j_k,
    \qquad u_k=hs_k.
\end{equation}
For the same left-rectangle convention, $j_k\geq1$ and $s_k\geq0$.
The discrete simplex integral is therefore
\begin{equation}
    \boxed{\widehat Q_n[m]\simeq h^n
    \sum_{j_1,\ldots,j_n=1}^{N-1}
    \sum_{s_1,\ldots,s_n=0}^{N-1}
    f_n[j_1,\ldots,j_n]
    \prod_{k=1}^n\Lambda[s_k]\,
        \delta_{s_{k-1},\,j_k+s_k},\qquad s_0=m.}
    \label{eq:quantics_ttn_simplex}
\end{equation}
The deltas enforce $\sum_k j_k\leq m$ and all intermediate ordered times.
For $\Lambda=1$, they impose only the simplex geometry; otherwise they
also supply the correct activation factors at the partial remaining times.
The sums in this formula specify a tensor-network contraction, not dense
loops over the $N$ possible values of each index.

Each delta is a binary addition tensor. At bit $\ell$ its local constraint is
\begin{equation}
    j_{k,\ell}+s_{k,\ell}+c_{k,\ell}
       =s_{k-1,\ell}+2c_{k,\ell+1},
    \qquad c_{k,0}=c_{k,B}=0,
\end{equation}
with carries $c_{k,\ell}\in\{0,1\}$. Zero boundary carries exclude
negative remaining times and overflow. Attach these addition tensors and
the MPSs for $\Lambda$ to the gap TTN, fix the bits of $s_0=m$, and
contract the other physical and virtual indices. Leaving $s_0$ open
instead defines a quantics output function, whose compression is a further
numerical task.

These constraints connect different branches through the remaining-time
indices. They are therefore \emph{not} the independent branch convolutions
used for rectangular integration: the augmented network is not generally
the original tree, and its contraction and compression need a separate
algorithm. Equation~\eqref{eq:quantics_ttn_simplex} describes that
construction; it is not an already implemented TTN operation.
If the saved TTN uses individual delays rather than gaps, the coordinate
transformation of Section~2 must also be accounted for; relabelling the
branches alone is not a change of variables.

\begin{thebibliography}{9}
\raggedright%
\bibitem{ttd}
Y.~N\'u\~nez-Fern\'andez et al.,
\emph{Learning Feynman Diagrams with Tensor Trains},
\href{https://arxiv.org/abs/2207.06135v2}{arXiv:2207.06135v2} (2022),
Sec.~V.A and Appendix G, Eqs.~(G1) to (G7);
see also Eqs.~(53) to (55).

\bibitem{jeannin}
M.~Jeannin,
\emph{Exact numerical calculation of Feynman diagrams for
out-of-equilibrium Kondo physics using Tensor Cross-Interpolation},
PhD thesis, Universit\'e Grenoble Alpes (2024),
Sec.~5.1.2, pp.~78--80, Eqs.~(5.10) to (5.17).
\end{thebibliography}

\end{document}
